\begin{definition}[Total complex phase]\label{mod:basic-phase}
For \(z\in\mathbf C\), define the \emph{total phase}
\[
 \operatorname{phase}(z)=
 \begin{cases}
  1,&z=0,\\[2mm]
  \dfrac{z}{\lvert z\rvert},&z\ne0.
 \end{cases}
\]
This is the phase operation used later in the definition of Langlands's
local constant.  For every \(z\in\mathbf C\),
\[
 \operatorname{phase}(z)\ne0,
 \qquad
 \bigl\lvert\operatorname{phase}(z)\bigr\rvert=1.
\]
If \(\lvert z\rvert=1\), then
\[
 \operatorname{phase}(z)=z.
\]
For \(z,w\ne0\),
\[
 \operatorname{phase}(zw)
 =\operatorname{phase}(z)\operatorname{phase}(w).
\]
The nonzero hypotheses are essential here because the phase has been
totalized by setting \(\operatorname{phase}(0)=1\); multiplicativity is
not asserted unconditionally.  For every \(n\in\mathbf N\),
\[
 \operatorname{phase}(z^n)=\operatorname{phase}(z)^n,
\]
and complex conjugation satisfies
\[
 \operatorname{phase}(\overline z)
 =\overline{\operatorname{phase}(z)}.
\]
For every positive real number \(r>0\), regarded as a complex number,
\[
 \operatorname{phase}(r)=1,
 \qquad
 \operatorname{phase}(rz)=\operatorname{phase}(z).
\]
Finally, if \(\lvert u\rvert=1\) and \(z\ne0\), then
\[
 \operatorname{phase}(uz)=u\,\operatorname{phase}(z).
\]
\LeanFile{LanglandsFirstMainLemma/Basic/Phase.lean}
\lean{LanglandsFirstMainLemma.phase}
\leanok
\SourceText{Local notation, subsection Definition of Delta.}
\uses{mod:prelude}
\FormalizationContract
\end{definition}
